\documentclass[10pt,a4paper]{amsart}
\usepackage[margin=19mm]{geometry}
\usepackage{amsmath,amssymb,mathtools}
\usepackage[scr=rsfs,cal=euler]{mathalfa}
\usepackage{xfrac}
\usepackage[unicode,hidelinks,bookmarks=false]{hyperref}
\newcommand{\ie}{i.e.\ }
\newcommand{\cf}{cf.\ }
\newcommand{\resp}{respectively\ }
\newcommand{\Subsection}{\subsection}
\providecommand{\nobreakdash}{\nobreak\hbox{-}}
\setlength{\emergencystretch}{2em}
\multlinegap0pt
\usepackage[shortlabels]{enumitem}
\usepackage{amssymb}
\usepackage{mathtools}
\usepackage{xcolor}
\newtheorem{lemma}{Lemma}
\newcommand{\R}{\ensuremath{\mathbb{R}}}
\title{Unified theory for regularity persistence of vortex patch boundaries}
\author{Marc Maga\~na}
\date{Source: arXiv:2608.02033v1 (CC BY 4.0)}
\begin{document}
\maketitle

\noindent\emph{Source and licence.} Unified theory for regularity persistence of vortex patch boundaries. Version arXiv:2608.02033v1 (CC BY 4.0), distributed by arXiv under the Creative Commons Attribution 4.0 International licence, http://creativecommons.org/licenses/by/4.0/, as recorded in the arXiv licence metadata for this exact version and shown on the abstract page https://arxiv.org/abs/2608.02033v1. Full licence notice: \texttt{licenses/arxiv-2608.02033-CC-BY-4.0.txt}.\par\smallskip

\noindent\textit{Editorial scope.} Counted unit: the article's lemma lemma:canviVar (source0.tex lines 510--516). The article's complete original proof of it is delivered with it (source0.tex lines 518--525, one \texttt{proof} environment of 8 lines). No mathematics is written, repaired or re-derived: the statement and the proof are the article's own lines, in the article's own notation, and no retained line is substituted or re-flowed.  Retained uncounted as the source-local context the proof needs: lines 172--177; lines 368--374.  Nothing was omitted from any retained span, so the package carries no omission record.  No retained line contains a citation, so the wrapper carries no bibliography and no bibliography entry.  No retained line contains a figure environment, an image-inclusion command or drawing code, so no image is delivered and the package declares no asset dependency.  Editorial formatting only, in addition: an AMS \texttt{amsart} preamble that restates the article's own declarations and macros used by the retained lines, namely the environments \texttt{lemma} on separate counters, together with the standard packages those lines need.  The retained environments print in this document as Lemma 1 in order of appearance, whereas the article prints the same items with the numbering of its own full document.  The complete native source of the article as distributed in the arXiv e-print for this exact version is bundled unchanged as \texttt{sources/206661/source0.tex}.
\begin{enumerate}[label={(A\arabic*)}]
    \item $K(r)\in C^\infty (\R\setminus \{0\}),$ \label{ass:smoothness}
    \item $\int_0^1 |K(r)| dr < \infty,$ \label{ass:integrability}
    \item There exists $ \delta >0$ such that $|K(r)|$ is decreasing on $(0,\delta]$. Moreover, for $r>\delta,$ there exists constants $C>0$ and $M\geq 0$ with $|K(r)|\leq C(1+r^M).$ \label{ass:compKernel}
    \item For $n=1,2,3$, $|K^{(n)}(r)|$ is decreasing for all $r>0$. \label{ass:monotonicity}
\end{enumerate}

\begin{lemma} \label{lemma:intDerivades}
    If $K$ satisfies \ref{ass:smoothness}--\ref{ass:monotonicity}, then
    \begin{enumerate}
        \item $\displaystyle \lim_{r\to 0} r^{n+1}K^{(n)}(r)=0$ for $n=0,1,2$, \label{prop:Limits0}
        \item $\displaystyle \int_0^1 r^n\,|K^{(n)}(r)|\,dr < \infty$ for $n=1,2,3$. \label{prop:integrabilitatrK}
    \end{enumerate}
\end{lemma}

\begin{lemma} \label{lemma:canviVar}
     Let $j\geq n$ with $n=1,2,3$. If $K$ satisfies \ref{ass:smoothness}-\ref{ass:monotonicity} then for all $S>0$ we have 
     \begin{equation*}
         \int_0^S \nu^j |K^{(n)}(\nu)|d\nu \leq \int_0^1 \nu^j |K^{(n)}(\nu)| d\nu + \frac{|K^{(n)}(1)|}{j+1}(S^{j+1}-1)_+ \leq C_1 + C_2(S^{j+1}-1)_+,
     \end{equation*}
     where $(x)_+=\max\{x,0\},$ and the constants $C_1, C_2$ depend on $K$, $n$ and $j$.
\end{lemma}

\begin{proof}
    If $S\leq 1$ then $(S^{j+1}-1)_+=0$ so the result is clear. For $S\geq 1$ we split the integral into two parts and use the monotonicity assumption \ref{ass:monotonicity} to get
    \begin{align*}
        \int_0^S \nu^j |K^{(n)}(\nu)|d\nu &\leq \int_0^1 \nu^j |K^{(n)}(\nu)| d\nu + |K^{(n)}(1)| \int_1^S \nu^j  d\nu \\
        & \leq \int_0^1 \nu^j |K^{(n)}(\nu)| d\nu + \frac{|K^{(n)}(1)|}{j+1}(S^{j+1}-1).
    \end{align*}
    Finally, by Lemma \ref{lemma:intDerivades} we have that the first term is bounded.
\end{proof}

\end{document}
